\section{Tool Calls as Executive Acts}
\label{sec:executive-acts}

A tool call is an executive act in a precise sense. The agent, holding delegated authority, dispatches that authority against a resource it does not own, producing an effect that, if wrong, must be undone.

\paragraph{The five-part structure.}
An executive act has five components. (1) An \emph{authorisation chain}: who delegated authority to the acting party, under what scope, with what limitations. (2) A \emph{subject}: the resource the act will affect. (3) An \emph{effect}: a state transition or external communication. (4) A \emph{claimed lifetime}: the duration the act is intended to remain in force, after which it is expected to be discharged. (5) A \emph{reversal expectation}: an implicit or explicit promise that the act can be undone if it should not have happened.

Tool calls instantiate each component. The authorisation chain is the credential and system prompt under which the agent runs, optionally enriched by a capability token. The subject is the tool server (database, filesystem, email recipient, payments endpoint). The effect is whatever state the call mutates. The lifetime is, in current deployments, implicit and infinite: a write is permanent unless explicitly undone. The reversal expectation is implicit and often false: agent platforms log tool calls but rarely capture the inverse operation that would undo them.

\paragraph{The structural parallel.}
The same five components appear in domains that have been formally analysed in adjacent literature. An emergency response action in an endpoint-detection-and-response (EDR) system has an authorisation chain (the security policy plus executor key), a subject (a process tree, a file, a network flow), an effect (quarantine, suspend, restrict egress), a claimed lifetime (a TTL after which the action is reviewed or withdrawn), and a reversal expectation (an inverse executor releasing the quarantine or resuming the process). A temporally-bounded administrative order in a legal system has the same shape: a delegated authority, a subject, an effect, a sunset clause, a reversal mechanism. An agent's call to delete a row from a customer database has the same structure.

The parallel is load-bearing rather than metaphorical. A formal grammar for bounded executive action has been developed for the EDR setting~\cite{programmableSovereignty2026}, and the grammar is the same for the tool-call case. The bounded-executive-action construction captures all three under one typed envelope.

\paragraph{The current deployment surface.}
Current tool-call stacks such as the Model Context Protocol~\cite{anthropicMCP2024} are passive dispatchers: schema-described function calls are emitted by the agent, dispatched by the runtime, and executed by the tool server. Deployments may add operator approval, a denylist, or per-tool rate limits, but they do not in general require the caller to construct the call with a typed rollback witness, do not enforce a positive TTL by construction, and do not produce a tamper-evident receipt that an external party can replay.

Treating the call as an executive act exposes guarantees the agent-as-policy framing obscures. A passive dispatcher can refuse calls only on properties of the agent that produced them (identity, rate, claimed prompt). A substrate that models the call as an executive act and demands its structural witnesses at admission can refuse calls that violate the grammar regardless of who produced them and regardless of whether the operator approved them.
